\section{Analytic and weighted propagation estimates}\label{sec:analytic}
Proposition~\ref{prop:dd-id}(ii) reduces the approximation error to a simplex integral of
$\|\mathcal V_{s_1}b'\|$, where $b'$ is a polynomial in $X,Y$ applied to the amplitude.
The estimates in this section serve the two terms of the comparison of $\mathcal V_{s_1}b'$ with the mollified flow
$V_\varepsilon$ made in the proofs of Section~\ref{sec:companion}: the term $\mathcal V_0V_\varepsilon(s_1)^*b'$, and a Duhamel
integral of the mollifier defect $(\Delta(r)-\Delta_\varepsilon(r))V_\varepsilon(r,s_1)b'$, $0\le r\le s_1$.
Weighted propagation (Proposition~\ref{prop:ws}) and periodization control the first term.
For the second, analytic propagation (Proposition~\ref{prop:a0}) and the mollifier-defect
estimate (Lemma~\ref{lem:unw}) give exponential smallness without a spatial weight;
Lemma~\ref{lem:path} gives a weighted bound. Interpolation between these bounds compensates
for the growth caused by periodization and transport. The derivative estimates of
Sections~\ref{subsec:dd-fc}, \ref{app:m} and~\ref{subsec:dd-wb} control $b'$ with the required
factorial dependence on the expansion and derivative orders.

\medskip\noindent\textbf{Notation.}
We use $X$, $Y$, $\Delta(s)$, $\Lambda$, $X_0$, $\tilde\psi$ and the power-weight norm $\pwNorm a=\norm{\ue^{\alpha\Lambda}X_0^\beta a}$
of Section~\ref{sec:setup}.
By Lemma~\ref{lem:m2w} a finite norm with the weight $\ue^{\beta\langle d\rangle}$ at radius $\alpha$ gives a finite norm of this
kind at the radius $c_0\alpha$; this gives the reversed-weight consequence in Section~\ref{sec:swapped}.
The two-time mollified propagator $V_\varepsilon(s_1,s_0)=V_\varepsilon(s_1)V_\varepsilon(s_0)^*$ is defined in
Section~\ref{app:kl}.

\subsection{The mollifier and elliptic regularity}
Let $J_\varepsilon:=j_\varepsilon(\Lambda)$ with $j_\varepsilon:=g*\chi(\varepsilon\,\cdot)$, where $g$ is the unit Gaussian and, as in
\cite[(4.20)--(4.21)]{S}, $\chi\in C_c^\infty(\mathbb R)$ has $\operatorname{supp}\chi\subset[-2,2]$ and $\chi=1$ on $[-1,1]$; in
addition $\chi$ is real, even and $0\le\chi\le1$, so $0\le j_\varepsilon\le1$. (In the notation of \cite{S},
$J_\varepsilon=\chi_\varepsilon(\varepsilon\Lambda)$ with $\chi_\varepsilon=g_\varepsilon*\chi$ and $g_\varepsilon=\varepsilon^{-1}g(\cdot/\varepsilon)$.) Put $\Delta_\varepsilon(t)=J_\varepsilon\Delta(t)J_\varepsilon$, with $\Delta(t)=\ue^{-2t}X^2+Y^2-Y$.

\begin{lemma}[elliptic estimate]\label{lem:ell}
For $f,g\in C_c^\infty(\Hb)$,
$\langle Xf,g\rangle=-\langle f,Xg\rangle$ and $\langle Yf,g\rangle=-\langle f,Yg\rangle+\langle f,g\rangle$, and
\[
\norm{X^2f}^2+\norm{(Y^2-Y)f}^2+2\norm{XYf}^2=\norm{\Delta f}^2+2\norm{Xf}^2\le\norm{\Delta f}^2+2\norm f\norm{\Delta f},
\]
and $\norm{Xf}^2+\norm{Yf}^2=\langle f,-\Delta f\rangle$. By closure, these three relations hold for every
$f\in\operatorname{Dom}(\Delta)$, and then $\norm{X^2f}$, $\norm{(Y^2-Y)f}\le\norm{\Delta f}+\norm f$,
$\norm{Y^2f}\le2(\norm{\Delta f}+\norm f)$, and
$\norm{\Delta(t)f}\le2(\norm{\Delta f}+\norm f)$ for $t\ge0$. Moreover
$\Delta_\varepsilon(t)$ is a bounded self-adjoint operator on $L^2(\Hb)$, norm continuous in $t\ge0$.
\end{lemma}
\begin{proof}
Let $f,g\in C_c^\infty(\Hb)$. Integration by parts with $d\nu=y^{-2}dx\,dy$ gives the two formulas for $X=y\partial_x$ and $Y=y\partial_y$;
that is, $X^*=-X$ and $Y^*=1-Y$ on test functions. So $-\Delta=X^*X+Y^*Y$, which gives
$\norm{Xf}^2+\norm{Yf}^2=\langle f,-\Delta f\rangle$, and
$Y^2-Y=-Y^*Y$. Since $[Y,X]=X$, $YX^2=X^2Y+2X^2$, and therefore
$\langle X^2f,Y^*Yf\rangle=\langle YX^2f,Yf\rangle=-\norm{XYf}^2-2\langle Xf,XYf\rangle$. Moreover $XY=YX-X$ and $Y+Y^*=1$ give
$\mathrm{Re}\langle Xf,XYf\rangle=\mathrm{Re}\langle Xf,YXf\rangle-\norm{Xf}^2=-\tfrac12\norm{Xf}^2$. With $\Delta=X^2-Y^*Y$,
\begin{align*}
\norm{\Delta f}^2&=\norm{X^2f}^2+\norm{Y^*Yf}^2+2\norm{XYf}^2+4\,\mathrm{Re}\langle Xf,XYf\rangle\\
&=\norm{X^2f}^2+\norm{(Y^2-Y)f}^2+2\norm{XYf}^2-2\norm{Xf}^2,
\end{align*}
which is the identity; the inequality holds because $\norm{Xf}^2\le\langle f,-\Delta f\rangle\le\norm f\norm{\Delta f}$. Now let
$f\in\operatorname{Dom}(\Delta)$. Since $\Delta$ is essentially self-adjoint on $C_c^\infty(\Hb)$, there are test functions $f_n\to f$ with
$\Delta f_n\to\Delta f$; by the inequality applied to $f_n-f_m$, the sequences $Xf_n$, $Yf_n$, $X^2f_n$, $(Y^2-Y)f_n$, $XYf_n$ converge in
$L^2$, to the distributional derivatives of $f$, and the identities and the inequality pass to the limit. The right side is at most
$(\norm{\Delta f}+\norm f)^2$, which gives the bounds for $X^2f$ and $(Y^2-Y)f$. Since $\norm{Yf}^2\le\norm f\norm{\Delta f}$,
$\norm{Yf}\le\tfrac12(\norm{\Delta f}+\norm f)$, which gives the bound for $Y^2f=(Y^2-Y)f+Yf$; and $\Delta(t)=\ue^{-2t}X^2+(Y^2-Y)$.
Finally $J_\varepsilon$ maps $L^2$ into $\operatorname{Dom}(\Delta)$ with
$\norm{\Delta J_\varepsilon}\le\sup_\lambda\lambda^2j_\varepsilon(\lambda)<\infty$, the operators $X^2$ and $Y^2-Y=-Y^*Y$ are symmetric on
$\operatorname{Dom}(\Delta)$, and $\Delta_\varepsilon(t)-\Delta_\varepsilon(t')=(\ue^{-2t}-\ue^{-2t'})J_\varepsilon X^2J_\varepsilon$.
\end{proof}

\subsection{Cauchy estimates for derivatives}\label{subsec:dd-fc}

We first control a single vector field with a loss of analytic radius, iterate the estimate for words,
and then include the power weight by the Leibniz rule.

\begin{lemma}[first-order Cauchy estimate]\label{lem:fc}
For $\tau_0<\pi/2$ there is $C_1=C_1(\tau_0)$ such that for $Z\in\{X,Y\}$ and $0\le\alpha''<\alpha'\le\tau_0$,
\[
\norm{\ue^{\alpha''\Lambda}Z\,\ue^{-\alpha'\Lambda}}_{L^2(\Hb)}\le\frac{C_1}{\alpha'-\alpha''}.
\]
\end{lemma}

\begin{proof}
The Kontorovich--Lebedev kernel calculation is given in Appendix~\ref{app:special}, Section~\ref{app:fc-proof}.
\end{proof}

\begin{lemma}[Cauchy estimate for words]\label{lem:wc}
For a word $Z_1\cdots Z_m$ with $Z_i\in\{X,Y\}$, and $0\le\alpha''<\alpha\le\tau_0<\pi/2$:
\[
\norm{\ue^{\alpha''\Lambda}Z_1\cdots Z_m\ue^{-\alpha\Lambda}}\le\Bigl(\frac{C_1m}{\alpha-\alpha''}\Bigr)^m\le\Bigl(\frac{\ue C_1}{\alpha-\alpha''}\Bigr)^m m!.
\]
\end{lemma}

\begin{proof}
Take equally spaced radii $\alpha''=r_0<r_1<\cdots<r_m=\alpha$, and insert $\ue^{-r_i\Lambda}\ue^{r_i\Lambda}$.
Each factor $\ue^{r_{i-1}\Lambda}Z_i\ue^{-r_i\Lambda}$ is bounded by $C_1m/(\alpha-\alpha'')$ (Lemma~\ref{lem:fc}). The identity extends from the spectral core of Appendix~\ref{app:fc-proof} by closure, since each $\ue^{r_i\Lambda}$ is closed. Finally $m^m\le\ue^m m!$. \qedhere
\end{proof}

Weights enter through the following algebra. Let $\mathfrak A:=\operatorname{span}\{y,\;x,\;(x^2+1)/y\}$.
\begin{itemize}
\item $\mathfrak A$ contains $X_0=\tfrac12\bigl((x^2+1)/y+y\bigr)$.
\item $\mathfrak A$ is invariant under $X$ and $Y$: $Xy=0$, $Xx=y$, $X((x^2+1)/y)=2x$; $Yy=y$, $Yx=0$, $Y((x^2+1)/y)=-(x^2+1)/y$.
\item In this basis the $\ell^1$-operator norms are $\norm{X}\le2$ and $\norm{Y}\le1$.
\item Since $X_0\ge y/2$, $X_0\ge|x|$ (because $x^2+y^2+1-2y|x|=(|x|-y)^2+1$), and $X_0\ge(x^2+1)/(2y)$, every $v=c_1y+c_2x+c_3(x^2+1)/y$ satisfies $|v|\le2\norm{c}_1X_0$.
\item Hence $|Z_BX_0|\le2\cdot2^{|B|}X_0$ for every word $Z_B$.
\end{itemize}

\begin{lemma}[Faà di Bruno bound]\label{lem:fdb}
For $\beta\ge0$ and a word $Z_S$ of length $j$:
\[
|X_0^\beta Z_SX_0^{-\beta}|\le B_\beta\,j!\,\rho_*^{-j},\qquad \rho_*:=\tfrac14\log(3/2)
\]
(any fixed value below $\tfrac12\log(3/2)$ works; $\rho_*$ does not depend on $\beta$), where $B_\beta$ depends only on $\beta$.
\end{lemma}

\begin{proof}
Expand by induction on $j$: $Z_S\varphi(X_0)=\sum_{\pi}\varphi^{(|\pi|)}(X_0)\prod_{B\in\pi}Z_BX_0$, the sum over partitions $\pi$ of $S$, where each block $B$ inherits the order of $S$.
With $\varphi(u)=u^{-\beta}$, $|\varphi^{(k)}(u)|=(\beta)_ku^{-\beta-k}$. So the bound is $2^j\sum_kS(j,k)(\beta)_k2^k$, with the Stirling numbers $S(j,k)$.
The exponential generating function of $\sum_kS(j,k)(\beta)_k2^k$ is $(3-2\ue^z)^{-\beta}$, which is holomorphic on $|z|<\log(3/2)$.
Cauchy's estimate on $|z|=\rho<\log(3/2)$ gives the bound $C_{\beta,\rho}\,j!\,\rho^{-j}$. Combined with the factor $2^j$, this gives $\rho_*=\rho/2$. \qedhere
\end{proof}

\begin{lemma}[weighted Cauchy estimate for words]\label{lem:sc}
For $\beta\ge0$, a word $Z_1\cdots Z_m$ and $0<\alpha\le\tau_0$, with $F:=X_0^\beta a$ and $\norm{\ue^{\alpha\Lambda}F}<\infty$:
\[
\norm{X_0^\beta Z_1\cdots Z_m a}\le B_\beta(m+1)\,m!\,(C_3/\alpha)^m\,\norm{\ue^{\alpha\Lambda}X_0^\beta a},
\qquad C_3:=\max(\ue C_1,\tau_0/\rho_*).
\]
\end{lemma}

\begin{proof}
By the Leibniz rule for words of derivations,
\[X_0^\beta Z_1\cdots Z_m(X_0^{-\beta}F)=\sum_{S\subseteq[m]}(X_0^\beta Z_SX_0^{-\beta})\,(Z_{S^c}F)\]
pointwise. Bound each term by $B_\beta|S|!\rho_*^{-|S|}\norm{Z_{S^c}F}$, and the second factor by Lemma~\ref{lem:wc} with $\alpha''=0$: $\norm{Z_{S^c}F}\le(\ue C_1/\alpha)^{m-|S|}(m-|S|)!\,\norm{\ue^{\alpha\Lambda}F}$.
Since $\sum_{|S|=j}j!(m-j)!=m!$, the sum is $\le B_\beta\,m!\sum_j\rho_*^{-j}(\ue C_1/\alpha)^{m-j}\le B_\beta(m+1)m!(C_3/\alpha)^m$. \qedhere
\end{proof}

\subsection{Multiplication by spatial weights}\label{app:m}
\begin{lemma}[analytic vectors from word bounds]\label{lem:expw}
Let $0<\alpha\le1$, $L\ge1$ and $B,N_0\ge0$. Let $G\in L^2(\Hb)$ be such that every word $Z_1\cdots Z_m$, $Z_i\in\{X,Y\}$, $m\ge0$,
applied to $G$ in the sense of distributions lies in $L^2$, with
\[
\norm{Z_1\cdots Z_mG}\le B\,(m+1)\,m!\,(L/\alpha)^m\,N_0 .
\]
Then $G\in\operatorname{Dom}(\ue^{r\Lambda})$ and $\norm{\ue^{r\Lambda}G}\le C_*BN_0$ for $r=\alpha/(4L)$, with an absolute constant $C_*$.
\end{lemma}
\begin{proof}
$\Delta$ is essentially self-adjoint on $C_c^\infty(\Hb)$, so $G\in\operatorname{Dom}(\Delta^n)$ for every $n$, and
$(-\Delta)^n=(-X^2-Y^2+Y)^n$ is a sum of $3^n$ words of length at most $2n$ with coefficients $\pm1$. Since
$m\mapsto(m+1)m!(L/\alpha)^m$ is increasing,
$\norm{\Lambda^{2n}G}=\norm{(-\Delta)^nG}\le3^nB(2n+1)(2n)!(L/\alpha)^{2n}N_0$, and so
\[
\sum_{n\ge0}\frac{r^{2n}}{(2n)!}\norm{\Lambda^{2n}G}\le BN_0\sum_{n\ge0}(2n+1)\Bigl(\frac{3}{16}\Bigr)^n\le2BN_0 .
\]
The partial sums of $\sum_n(r\lambda)^{2n}/(2n)!$ increase to $\cosh(r\lambda)$, so monotone convergence in the spectral representation
of $\Lambda$ gives $\norm{\cosh(r\Lambda)G}\le2BN_0$, and $\ue^{r\lambda}\le2\cosh(r\lambda)$. So $C_*=4$ works.
\end{proof}

\begin{lemma}[removing a spatial weight with loss of analytic radius]\label{lem:mw}
Let $\beta\ge0$ and $\tau_0<\pi/2$, let $B_\beta$, $C_3$ be the constants of Lemma~\ref{lem:sc}, and put $c_M:=1/(4\max(C_3,1))$. For
$0<\alpha\le\min(1,\tau_0)$ and every $a$ with $\norm{\ue^{\alpha\Lambda}X_0^\beta a}<\infty$,
\[
\norm{\ue^{c_M\alpha\Lambda}a}\le C_*B_\beta\,\norm{\ue^{\alpha\Lambda}X_0^\beta a},
\]
with $C_*$ the constant of Lemma~\ref{lem:expw}.
\end{lemma}
\begin{proof}
By Lemma~\ref{lem:sc} and $X_0\ge1$, every word satisfies
$\norm{Z_1\cdots Z_ma}\le\norm{X_0^\beta Z_1\cdots Z_ma}\le B_\beta(m+1)m!(C_3/\alpha)^m\norm{\ue^{\alpha\Lambda}X_0^\beta a}$.
Apply Lemma~\ref{lem:expw} with $L=\max(C_3,1)$.
\end{proof}

\begin{lemma}[comparison of the spatial weights]\label{lem:m2w}
For $\beta\ge0$ there are $c_0\in(0,1]$ and $C$, depending only on $\beta$, such that for every base point $x_0$, every
$0<\alpha\le1$ and every $a$ with $\norm{\ue^{\alpha\Lambda}\ue^{\beta\langle d\rangle}a}<\infty$,
\[
\norm{\ue^{c_0\alpha\Lambda}X_0^\beta a}\le C\,\norm{\ue^{\alpha\Lambda}\ue^{\beta\langle d\rangle}a}.
\]
\end{lemma}
\begin{proof}
The fields $X,Y$ and the operator $\Lambda$ commute with the affine isometries $z\mapsto pz+q$, $p>0$, which act transitively on
$\Hb$; so we may take $x_0=i$. Put $F:=\ue^{\beta\langle d\rangle}a\in\operatorname{Dom}(\ue^{\alpha\Lambda})$. Then $X_0^\beta a=mF$ with
\[
m:=X_0^\beta\ue^{-\beta\langle d\rangle}=\Phi(\tilde\psi),\qquad\Phi(\psi):=\exp\beta\bigl(\psi-(1+d^2)^{1/2}+1\bigr),
\]
where $d\ge0$ is the function of $\psi\ge0$ given by $\psi=\log\cosh d$.

\emph{Step 1: $|\Phi^{(k)}|\le A\,k!\,\rho^{-k}$ on $[0,\infty)$, with $A,\rho>0$ depending only on $\beta$.} The function $\Phi$
extends holomorphically to a neighbourhood of every point of $[0,\infty)$. Near $0$ this holds because $d^2$ is the analytic inverse at
$0$ of $s\mapsto\log\cosh\sqrt s$, which has derivative $\tfrac12$ there. Near a point of $(0,\infty)$ it holds because
$d=\operatorname{arcosh}(\ue^\psi)$ is analytic there. On the strip $\{\operatorname{Re}\psi>\tfrac{17}{10},\ |\operatorname{Im}\psi|<\tfrac14\}$,
\begin{gather*}
\Phi=\exp\beta\bigl(1-\lambda-(D+R)^{-1}\bigr),\\
\lambda=\log\bigl(1+(1-\ue^{-2\psi})^{1/2}\bigr),\qquad D=\psi+\lambda,\qquad R=(1+D^2)^{1/2},
\end{gather*}
with principal branches: on the real axis $D=d$ and $\psi-R+1=1-\lambda-(D+R)^{-1}$, because $R^2-D^2=1$. On the strip
$\operatorname{Re}(1-\ue^{-2\psi})>0$, $\operatorname{Re}\lambda\ge0$, $\operatorname{Re}(1+D^2)>0$ and $\operatorname{Re}(D+R)>\tfrac{17}{10}$, so $\Phi$
is holomorphic there with $|\Phi|\le\ue^{2\beta}$. Compactness of $[0,2]$, and Cauchy's estimate on discs of radius $\tfrac15$ around
the points of $[2,\infty)$, give Step~1.

\emph{Step 2: $|Z_Sm|\le A\,c^j\,j!$ for every word $Z_S$ of length $j$, with $c=4+3/\rho$.} Put $\xi:=x/X_0$ and $\upsilon:=y/X_0$, so
$|\xi|\le1$, $0<\upsilon\le2$, $X\tilde\psi=\xi$ and $Y\tilde\psi=\upsilon-1$. The polynomials in $(\xi,\upsilon)$ are stable under $X$ and $Y$:
$X\xi=\upsilon-\xi^2$, $X\upsilon=-\xi\upsilon$, $Y\xi=\xi-\xi\upsilon$, $Y\upsilon=2\upsilon-\upsilon^2$. For
$Q=\sum c_{ab}\xi^a\upsilon^b$ let $\norm Q_*:=\sum|c_{ab}|2^b$, so $|Q|\le\norm Q_*$ pointwise. The four formulas give
$\norm{ZQ}_*\le4(\deg Q)\norm Q_*$, and $\norm{(Z\tilde\psi)Q}_*\le3\norm Q_*$. By induction on $j$,
$Z_S\Phi(\tilde\psi)=\sum_{k=0}^j\Phi^{(k)}(\tilde\psi)\,Q_{S,k}(\xi,\upsilon)$ with polynomials of degree at most $j$ satisfying
$Q_{ZS,k}=(Z\tilde\psi)Q_{S,k-1}+ZQ_{S,k}$. With $q_{j,k}:=\max_{|S|=j}\norm{Q_{S,k}}_*$ this gives
$q_{j+1,k}\le3q_{j,k-1}+4j\,q_{j,k}$, and then $T_j:=\sum_kk!\,\rho^{-k}q_{j,k}$ satisfies
$T_{j+1}\le(3(j+1)/\rho+4j)T_j\le c\,(j+1)\,T_j$, so $T_j\le c^jj!$. By Step~1, $|Z_Sm|\le A\,T_j$.

\emph{Step 3.} By the Leibniz rule for words, Step~2 and Lemma~\ref{lem:wc} at $\alpha''=0$ (with $\tau_0=1$), for a word of length $j$,
\[
\begin{aligned}
\norm{Z_S(mF)}&\le\sum_{S'\subseteq S}A\,c^{|S'|}|S'|!\,\Bigl(\frac{\ue C_1}{\alpha}\Bigr)^{j-|S'|}(j-|S'|)!\,\norm{\ue^{\alpha\Lambda}F}
\\ &\le A\,(j+1)\,j!\,(L/\alpha)^j\norm{\ue^{\alpha\Lambda}F},
\end{aligned}
\]
with $L:=\max(c,\ue C_1,1)$. Lemma~\ref{lem:expw} gives the claim with $c_0=1/(4L)$ and $C=C_*A$.
\end{proof}

\subsection{Analyticity of the mollified flow}\label{app:kl}
By Lemma~\ref{lem:ell} the equation $i\partial_sv=-\tfrac\hbar2\Delta_\varepsilon(s)v$ has a unitary propagator
$V_\varepsilon(s_1,s_0)$ on $L^2(\Hb)$, $s_0,s_1\ge0$, norm differentiable in both variables. The mollified flow of
\cite[(4.32)--(4.33)]{S} is $V_\varepsilon(s)=V_\varepsilon(s,0)$, and $V_\varepsilon(s_1,s_0)=V_\varepsilon(s_1)V_\varepsilon(s_0)^*$.
All estimates below are stated for $V_\varepsilon(s_1,s_0)$; this contains $V_\varepsilon(t)$ and $V_\varepsilon(t)^*$.

\medskip\noindent\textbf{The Kontorovich--Lebedev picture.}
Fourier transform in $x$; on the mode $\ue^{ikx}$, $k\ne0$, put $z=|k|y$ and write the function of $y$ as $z^{1/2}f(z)$. Then
$L^2(dx\,dy/y^2)$ becomes $L^2(dz/z)$ in $f$, $-\Delta$ becomes $-(z\partial_z)^2+z^2+\tfrac14$, and $X$ becomes multiplication by
$i\operatorname{sgn}(k)\,z$; on the mode $k=0$, $X=0$. The Kontorovich--Lebedev transform $F(\nu)=\int_0^\infty K_{i\nu}(z)f(z)\,dz/z$, with Parseval
$\int|f|^2dz/z=\int|F|^2w(\nu)\,d\nu$, $w(\nu)=\tfrac2{\pi^2}\nu\sinh(\pi\nu)$ (\cite[eqs.~(1.1)--(1.3)]{Y}, applied to $f(z)/z$), diagonalizes
$-\Delta=\nu^2+\tfrac14$. So $\Lambda$ acts as multiplication by $\Lambda(\nu)=(\nu^2+\tfrac14)^{1/2}$, and
$|\Lambda(\nu)-\Lambda(\nu')|\le|\nu-\nu'|$. In the unitary picture $G=w^{1/2}F$, multiplication by $z$ has the kernel
$k_1(\nu,\nu')$ computed in step (X) of the proof of Lemma~\ref{lem:fc}, and
\begin{equation}\label{eq:zt1}
0\le k_1(\nu,\nu')\le\frac{(\nu\nu')^{1/2}}{2\cosh(\pi(\nu-\nu')/2)} .
\end{equation}
The operators are the same for every $k$, so all bounds are uniform in $k$.

\begin{lemma}[clipped conjugation of $XJ_\varepsilon$]\label{lem:clip}
Let $\tau_0<\pi/2$, $L\ge\tfrac12$ and $\Lambda_{\le L}:=\min(\Lambda,L)$. For real $s$ let $R_s:=\ue^{s\Lambda_{\le L}}XJ_\varepsilon\ue^{-s\Lambda_{\le L}}$, a bounded
operator. Put
\begin{gather*}
g_K(x):=\frac{(1+|x|)\,\ue^{\tau_0|x|}}{2\cosh(\pi x/2)},\qquad g_Q(x):=\frac{(1+3|x|)\,|x|\,\ue^{\tau_0|x|}}{2\cosh(\pi x/2)},\\
A_1:=\int_{\mathbb R}g_K(x)\,dx,\qquad A_2:=\int_{\mathbb R}g_Q(x)\,dx .
\end{gather*}
For $0\le\tau\le\tau_0$, $\varepsilon>0$ and $u\in L^2(\Hb)$,
\[
\bigl\|\tfrac12(R_\tau+R_{-\tau})u\bigr\|\le A_1\norm{\Lambda J_\varepsilon u},\qquad
\bigl\|\tfrac12(R_\tau-R_{-\tau})u\bigr\|\le A_2\,\tau\,\norm{\Lambda_{\le L}J_\varepsilon u},
\]
and consequently
\[
\bigl|\mathrm{Im}\langle R_{-\tau}u,R_\tau u\rangle\bigr|\le2A_1A_2\,\tau\,\norm{\Lambda J_\varepsilon u}\,\norm{\Lambda_{\le L}J_\varepsilon u}.
\]
\end{lemma}
\begin{proof}
In the Kontorovich--Lebedev picture $\Lambda_{\le L}$ and $J_\varepsilon$ are multiplications by $\Lambda_{\le L}(\nu)=\min(\Lambda(\nu),L)$ and
$j_\varepsilon(\Lambda(\nu))$. Put $x=\nu-\nu'$ and $\delta=\Lambda_{\le L}(\nu)-\Lambda_{\le L}(\nu')$, so $|\delta|\le|x|$. Up to the factor
$i\operatorname{sgn}(k)$ of modulus one, the operators $\tfrac12(R_\tau+R_{-\tau})$ and $\tfrac12(R_\tau-R_{-\tau})$ have the kernels
\[
k_1(\nu,\nu')\cosh(\tau\delta)\,j_\varepsilon(\Lambda(\nu'))\qquad\text{and}\qquad k_1(\nu,\nu')\sinh(\tau\delta)\,j_\varepsilon(\Lambda(\nu')).
\]

\emph{Even part.} Since $(\nu\nu')^{1/2}\le\nu'+|x|/2$ and $\Lambda(\nu')\ge\max(\nu',\tfrac12)$, we have
$(\nu\nu')^{1/2}\le(1+|x|)\Lambda(\nu')$. By \eqref{eq:zt1} and $\cosh(\tau\delta)\le\ue^{\tau_0|x|}$, the even part is
$K_A\circ\Lambda J_\varepsilon$ with a kernel $|K_A(\nu,\nu')|\le g_K(\nu-\nu')$, and the Schur test gives $\norm{K_A}\le A_1$.

\emph{Odd part.} Here $|\sinh(\tau\delta)|\le\tau|\delta|\,\ue^{\tau_0|x|}$, and $\delta=0$ when $\Lambda(\nu)\ge L$ and
$\Lambda(\nu')\ge L$. In the remaining cases $(\nu\nu')^{1/2}\le(1+|x|)\,\Lambda_{\le L}(\nu')$: if $\Lambda(\nu')\le L$ this is the inequality
of the even part; if $\Lambda(\nu')>L\ge\Lambda(\nu)$, then $\Lambda_{\le L}(\nu')=L$ and
$(\nu\nu')^{1/2}\le\nu+|x|/2\le L+|x|L$, because $\nu\le\Lambda(\nu)\le L$ and $L\ge\tfrac12$. So the odd part is $K_B\circ\Lambda_{\le L}J_\varepsilon$ with
$|K_B(\nu,\nu')|\le\tau\,|x|\,g_K(x)\le\tau\,g_Q(x)$, and the Schur test gives $\norm{K_B}\le A_2\tau$.

\emph{The imaginary part.} With $P=\tfrac12(R_\tau+R_{-\tau})u$ and $Q=\tfrac12(R_\tau-R_{-\tau})u$,
\[\langle R_{-\tau}u,R_\tau u\rangle=\langle P-Q,P+Q\rangle=\norm P^2-\norm Q^2+2i\,\mathrm{Im}\langle P,Q\rangle,\]
so
$|\mathrm{Im}\langle R_{-\tau}u,R_\tau u\rangle|\le2\norm P\norm Q$.
\end{proof}

\begin{proposition}[analytic propagation for the mollified flow]\label{prop:a0}
Let $\tau_0<\pi/2$ and $c_A:=(1+32A_1A_2)^{-1}$, with $A_1,A_2$ as in Lemma~\ref{lem:clip}. Let $0<\varepsilon\le1$, $\hbar>0$,
$0\le\alpha'\le\alpha\le\tau_0$ with $\alpha>0$, and $s_0,s_1\ge0$ with
\[
|s_1-s_0|\le c_A\,\frac{\alpha-\alpha'}{\alpha}\,\frac{\varepsilon}{\hbar}.
\]
If $a\in\operatorname{Dom}(\ue^{\alpha\Lambda})$, then $V_\varepsilon(s_1,s_0)a\in\operatorname{Dom}(\ue^{\alpha'\Lambda})$ and
\[
\norm{\ue^{\alpha'\Lambda}V_\varepsilon(s_1,s_0)a}\le\norm{\ue^{\alpha\Lambda}a}.
\]
\end{proposition}
\begin{proof}
\emph{A frequency bound.} For $L\ge4/\varepsilon$ and $\lambda\ge0$, $\lambda^2j_\varepsilon(\lambda)^2\le(32/\varepsilon)\min(\lambda,L)$.
Indeed $j_\varepsilon(\lambda)\le\int_{|\mu|\le2/\varepsilon}g(\lambda-\mu)\,d\mu\le1$. If $\lambda\le4/\varepsilon$ the left side is at most
$(4/\varepsilon)\lambda$. If $\lambda>4/\varepsilon$, then $\lambda-\mu>\lambda/2$ on the range of integration, so
$j_\varepsilon(\lambda)\le\ue^{-\lambda^2/8}$ and $\lambda^2j_\varepsilon(\lambda)^2\le4/\ue$, while $\min(\lambda,L)\ge4/\varepsilon\ge4$.

\emph{The clipped energy.} Fix $L\ge4/\varepsilon$, put $\Lambda_{\le L}=\min(\Lambda,L)$ and $\rho:=32A_1A_2\hbar/\varepsilon$, and let
$\tau(s):=\alpha\,\ue^{-\rho|s-s_0|}$, $v(s):=V_\varepsilon(s,s_0)a$, $u(s):=\ue^{\tau(s)\Lambda_{\le L}}v(s)$. All operators involved are bounded.
Since $v'=\tfrac{i\hbar}2\Delta_\varepsilon(s)v$,
\[
\tfrac{d}{ds}\norm{u}^2=2\tau'\langle u,\Lambda_{\le L}u\rangle-\hbar\,\mathrm{Im}\bigl\langle u,\ue^{\tau\Lambda_{\le L}}\Delta_\varepsilon(s)\ue^{-\tau\Lambda_{\le L}}u\bigr\rangle .
\]
Now $\Delta_\varepsilon(s)=J_\varepsilon\Delta J_\varepsilon-(1-\ue^{-2s})J_\varepsilon X^2J_\varepsilon$. The first term commutes with $\Lambda_{\le L}$ and is
self-adjoint, so it does not contribute. For the second, $\ue^{\tau\Lambda_{\le L}}J_\varepsilon X^2J_\varepsilon\ue^{-\tau\Lambda_{\le L}}=-R_{-\tau}^*R_\tau$,
because $X^*=-X$. Hence, by Lemma~\ref{lem:clip}, $\norm{\Lambda_{\le L}J_\varepsilon u}\le\norm{\Lambda J_\varepsilon u}$ and the frequency bound,
for $s\ge0$,
\[
\begin{aligned}
\Bigl|\tfrac{d}{ds}\norm u^2-2\tau'\langle u,\Lambda_{\le L}u\rangle\Bigr|&=\hbar(1-\ue^{-2s})\bigl|\mathrm{Im}\langle R_{-\tau}u,R_\tau u\rangle\bigr|\\ &\le2A_1A_2\hbar\,\tau\,\langle u,\Lambda^2J_\varepsilon^2u\rangle\le2\rho\,\tau\,\langle u,\Lambda_{\le L}u\rangle .
\end{aligned}
\]
Since $\tau'=-\rho\tau$ for $s>s_0$ and $\tau'=\rho\tau$ for $s<s_0$, $\norm{u(s)}$ does not increase as $s$ moves away from $s_0$, in
either direction; all times between $s_0$ and $s_1$ are nonnegative. Therefore
$\norm{\ue^{\tau(s_1)\Lambda_{\le L}}v(s_1)}\le\norm{\ue^{\alpha\Lambda_{\le L}}a}\le\norm{\ue^{\alpha\Lambda}a}$.

\emph{The radius.} By the hypothesis on $|s_1-s_0|$ and $32A_1A_2c_A\le1$,
$\tau(s_1)\ge\alpha(1-\rho|s_1-s_0|)\ge\alpha'$. So $\norm{\ue^{\alpha'\min(\Lambda,L)}v(s_1)}\le\norm{\ue^{\alpha\Lambda}a}$ for every
$L\ge4/\varepsilon$, and monotone convergence in the spectral representation of $\Lambda$, as $L\to\infty$, gives the claim.
\end{proof}

\subsection{Weighted bounds for the mollifier and the mollified flow}\label{app:ws}
In this section the base point is $x_0=i$ (Lemma~\ref{lem:base} reduces the general case to it), $X_0=\cosh d(i,\cdot)$, and
$D:=(-\Delta-\tfrac14)^{1/2}$, so that $\Lambda=(D^2+\tfrac14)^{1/2}$. Explicitly $X_0=(x^2+y^2+1)/(2y)$, $XX_0=x$, $YX_0=y-X_0$, and
$(XX_0)^2+(YX_0)^2=X_0^2-1$.

\medskip\noindent\textbf{Bounded weights.} A \emph{bounded weight} is a smooth function $W>0$ on $\Hb$ such that $Z_SW$ is bounded
for every word $Z_S$ in $X,Y$, the empty word included. For $\rho,K\ge0$ let $\mathcal W(\rho,K)$ be the set of bounded weights
with, for $Z\in\{X,Y\}$,
\[
|ZW|\le\rho\,W,\qquad|Z^2W|\le K\,W,\qquad|\Delta W|\le K\,W,\qquad|Z^2(W^2)|\le K\,W^2 .
\]
For $\gamma\ge0$ and an integer $n\ge1$ put $W_{\gamma,n}:=(X_0^{-\gamma}+1/n)^{-1}$. Then $W_{\gamma,n}\le n$,
$W_{\gamma,n}\uparrow X_0^\gamma$ as $n\to\infty$,
\[
(X\log W_{\gamma,n})^2+(Y\log W_{\gamma,n})^2\le\gamma^2,
\]
and $W_{\gamma,n}\in\mathcal W(\gamma,K_\gamma)$ with a constant $K_\gamma$ depending only on $\gamma$ ($K_\gamma=8\gamma^2+7\gamma$ works).
Indeed, with $\phi=X_0^{-\gamma}$ one has $(X\phi)^2+(Y\phi)^2\le\gamma^2\phi^2$ and, since $X^2X_0=y\le2X_0$ and $Y^2X_0=X_0$,
$|Z^2\phi|\le(\gamma^2+3\gamma)\phi$; now differentiate $W_{\gamma,n}=(\phi+1/n)^{-1}$ and $W_{\gamma,n}^2$, and use $\phi\le\phi+1/n$.
All words applied to $\phi$ are bounded (Lemma~\ref{lem:fdb}), hence so are all words applied to $W_{\gamma,n}$.
A bound for all $W_{\gamma,n}$ with a constant independent of $n$ passes to the weight $X_0^\gamma$ by monotone convergence. The
bounds below are used in this way.

\begin{lemma}[wave energy at a bounded weight]\label{lem:wave}
Let $W$ be a bounded weight, $u\in L^2(\Hb)$ and $s\in\mathbb R$.
\begin{itemize}
\item[(a)] If $|XW|\le\rho W$ and $|YW|\le\rho W$, then $\norm{W\cos(s\Lambda)u}\le\ue^{2\rho|s|}\norm{Wu}$. If moreover $u$ lies in the
range of a spectral projection $1_{[0,R]}(\Lambda)$, $R<\infty$, then, with $q=\Lambda^{-1}\sin(s\Lambda)u$,
\[
\norm{WXq}\le\ue^{2\rho|s|}\norm{Wu},\qquad\norm{WYq}\le\ue^{2\rho|s|}\norm{Wu}.
\]
\item[(b)] If $l\ge0$ and $(X\log W)^2+(Y\log W)^2\le l^2$, then $\norm{W\cos(sD)u}\le\ue^{(l+1/2)|s|}\norm{Wu}$.
\end{itemize}
\end{lemma}
\begin{proof}
First let $u$ lie in the range of a spectral projection $1_{[0,L]}(\Lambda)$. Then all functions below are smooth, lie in the domain of
every power of $\Delta$, and are differentiable in $s$. For $f\in\operatorname{Dom}(\Delta)$ and $g$ in the form domain,
$\langle g,\Delta f\rangle=-\langle Xg,Xf\rangle-\langle Yg,Yf\rangle$ (Lemma~\ref{lem:ell}), and $W^2g$ is in the form domain when $g$ is,
since $W$ and $XW$, $YW$ are bounded.

(a) Let $p=\cos(s\Lambda)u$, so $q'=p$ and $p'=\Delta q$, and $E:=\norm{Wp}^2+\norm{WXq}^2+\norm{WYq}^2$. Then
\begin{align*}
E'&=2\,\mathrm{Re}\langle W^2p,\Delta q\rangle+2\,\mathrm{Re}\langle W^2Xq,Xp\rangle+2\,\mathrm{Re}\langle W^2Yq,Yp\rangle\\
&=-2\,\mathrm{Re}\langle X(W^2)p,Xq\rangle-2\,\mathrm{Re}\langle Y(W^2)p,Yq\rangle .
\end{align*}
Since $|X(W^2)|\le2\rho W^2$ and $|Y(W^2)|\le2\rho W^2$, $|E'|\le4\rho\norm{Wp}(\norm{WXq}+\norm{WYq})\le4\rho E$, so
$E(s)\le\ue^{4\rho|s|}E(0)=\ue^{4\rho|s|}\norm{Wu}^2$.

(b) Let $p=\cos(sD)u$ and $q=\int_0^s\cos(rD)u\,dr$, so $q'=p$ and $p'=(\Delta+\tfrac14)q$, and
$E:=\norm{Wp}^2+\norm{WXq}^2+\norm{WYq}^2+\tfrac14\norm{Wq}^2$. The same computation gives
\[E'=-2\,\mathrm{Re}\langle X(W^2)p,Xq\rangle-2\,\mathrm{Re}\langle Y(W^2)p,Yq\rangle+\mathrm{Re}\langle W^2p,q\rangle.\]
Pointwise
$|X(W^2)\,Xq+Y(W^2)\,Yq|\le2l\,W^2(|Xq|^2+|Yq|^2)^{1/2}$, so the first two terms are at most
$4l\norm{Wp}(\norm{WXq}^2+\norm{WYq}^2)^{1/2}\le2lE$ together, and $|\langle W^2p,q\rangle|\le\norm{Wp}\norm{Wq}\le E$. So $|E'|\le(2l+1)E$.

For general $u$ the bounds for $\cos(s\Lambda)u$ and $\cos(sD)u$ follow by applying this to $u_L=1_{[0,L]}(\Lambda)u$ and letting
$L\to\infty$, since $W$ is bounded.
\end{proof}

\begin{lemma}[commutator with the wave group]\label{lem:wcomm}
Let $W\in\mathcal W(\rho,K)$. For $u,g$ in the range of a spectral projection $1_{[0,R]}(\Lambda)$, $R<\infty$, and $s\in\mathbb R$,
\[
\bigl|\langle u,[\cos(s\Lambda),W^2]g\rangle\bigr|\le\bigl(8\rho|s|+(2K+2\rho)s^2\bigr)\ue^{4\rho|s|}\norm{Wu}\norm{Wg}.
\]
\end{lemma}
\begin{proof}
Put $h=W^2$, $c(s)=\cos(s\Lambda)$ and $Q(s)=\Lambda^{-1}\sin(s\Lambda)$. The function
$v(s)=[c(s),h]g$ satisfies $v''=\Delta v+[\Delta,h]c(s)g$, $v(0)=v'(0)=0$, so
\[
v(s)=\int_0^sQ(s-r)[\Delta,h]c(r)g\,dr.
\]
Here
$[\Delta,h]f=(\Delta h)f+2(Xh)Xf+2(Yh)Yf$, and $[\Delta,h]^*=-[\Delta,h]$. With $q(r):=Q(r)u$,
\[
\langle u,v(s)\rangle=-\int_0^s\Bigl(\bigl\langle(\Delta h)\,q(s-r),c(r)g\bigr\rangle+2\bigl\langle(Xh)Xq(s-r)+(Yh)Yq(s-r),c(r)g\bigr\rangle\Bigr)dr,
\]
in which no derivative falls on $c(r)g$. Now
$|Xh|,|Yh|\le2\rho h$ and $|\Delta h|\le|X^2h|+|Y^2h|+|Yh|\le(2K+2\rho)h$. By Lemma~\ref{lem:wave}(a),
$\norm{Wc(r)g}\le\ue^{2\rho|r|}\norm{Wg}$, $\norm{WXq(r)}+\norm{WYq(r)}\le2\ue^{2\rho|r|}\norm{Wu}$, and
$\norm{Wq(r)}\le|r|\ue^{2\rho|r|}\norm{Wu}$ because $q(r)=\int_0^rc(r')u\,dr'$. So the integrand is at most
$\bigl((2K+2\rho)|s-r|+8\rho\bigr)\ue^{2\rho|s|}\norm{Wu}\norm{Wg}$. Integrating over $r$ gives the bound
$\bigl(8\rho|s|+(K+\rho)s^2\bigr)\ue^{2\rho|s|}\norm{Wu}\norm{Wg}$, which implies the claim.
\end{proof}

By Fourier inversion, since $j_\varepsilon$ is even,
\[
J_\varepsilon=\frac1\pi\int_0^\infty h_\varepsilon(s)\cos(s\Lambda)\,ds,\qquad h_\varepsilon(s)=\varepsilon^{-1}\hat\chi(s/\varepsilon)\,\ue^{-s^2/2},
\]
where $h_\varepsilon$ is real, even and Schwartz. For $0<\varepsilon\le1$, $k,r\ge0$ and $C\ge0$,
\begin{equation}\label{eq:hmom}
\int_0^\infty|h_\varepsilon^{(r)}(s)|\,\ue^{Cs}\,ds\le C_{r}\,\varepsilon^{-r},\qquad\int_0^\infty|h_\varepsilon(s)|\,s^k\,\ue^{Cs}\,ds\le C_{k}\,\varepsilon^{k},
\end{equation}
with constants depending on $C$ and $\chi$ (substitute $s=\varepsilon\sigma$ and use $\ue^{Cs-s^2/2}\le\ue^{C^2/2}$).

\begin{lemma}[weighted mollifier]\label{lem:wmoll}
Let $W\in\mathcal W(\rho,K)$ and $0<\varepsilon\le1$. There is $C_W=C_W(\rho,K)$ such that for all $u,v,g\in L^2(\Hb)$ and
$Z\in\{X,Y\}$:
\begin{itemize}
\item[(i)] $\norm{WJ_\varepsilon u}\le C_W\norm{Wu}$;
\item[(ii)] $\norm{WZJ_\varepsilon u}\le C_W\varepsilon^{-1}\norm{Wu}$;
\item[(iii)] $\norm{W\Delta J_\varepsilon u}\le C_W\varepsilon^{-2}\norm{Wu}$ and $\norm{WZ^2J_\varepsilon u}\le C_W\varepsilon^{-2}\norm{Wu}$;
\item[(iv)] $|\langle v,[J_\varepsilon,W^2]g\rangle|\le C_W\varepsilon\norm{Wv}\norm{Wg}$.
\end{itemize}
In particular $\norm{W\Delta(s)J_\varepsilon u}\le3C_W\varepsilon^{-2}\norm{Wu}$ for $s\ge0$.
\end{lemma}
\begin{proof}
In (ii) and (iv) it suffices to take $u$, $v$, $g$ in the range of a spectral projection $1_{[0,R]}(\Lambda)$, with $C_W$ independent of
$R$: the projection commutes with $J_\varepsilon$, the operators $ZJ_\varepsilon$ and $[J_\varepsilon,W^2]$ are bounded on $L^2$, and $W$ is bounded.
(i) follows from Lemma~\ref{lem:wave}(a) and \eqref{eq:hmom}. (ii) With $q(s)=\Lambda^{-1}\sin(s\Lambda)u$, $q'=\cos(s\Lambda)u$ and
$q(0)=0$, an integration by parts gives $ZJ_\varepsilon u=-\tfrac1\pi\int_0^\infty h_\varepsilon'(s)\,Zq(s)\,ds$; use
Lemma~\ref{lem:wave}(a) and \eqref{eq:hmom} with $r=1$. (iii) Since $\Delta\cos(s\Lambda)u=\partial_s^2\cos(s\Lambda)u$ and
$h_\varepsilon'(0)=0$, two integrations by parts give $\Delta J_\varepsilon u=\tfrac1\pi\int_0^\infty h_\varepsilon''(s)\cos(s\Lambda)u\,ds$, and
\eqref{eq:hmom} with $r=2$ gives the first bound. For the second put $f=J_\varepsilon u$. The function $Wf$ lies in
$\operatorname{Dom}(\Delta)$, and
\begin{align*}
Z^2(Wf)&=WZ^2f+2(ZW)Zf+(Z^2W)f,\\
\Delta(Wf)&=W\Delta f+2(XW)Xf+2(YW)Yf+(\Delta W)f,
\end{align*}
with $|ZW|\le\rho W$, $|Z^2W|\le KW$ and $|\Delta W|\le KW$. So Lemma~\ref{lem:ell}, applied to $Wf$, bounds
$\norm{WZ^2f}$ by a constant times $\norm{W\Delta f}+\max_Z\norm{WZf}+\norm{Wf}$, and (i), (ii) and the first bound of (iii) apply. (iv) follows from
Lemma~\ref{lem:wcomm} and \eqref{eq:hmom} with $k=1,2$. The last claim holds because $\Delta(s)=\ue^{-2s}X^2+Y^2-Y$.
\end{proof}

\begin{proposition}[weighted propagation for the mollified flow]\label{prop:ws}
Let $W\in\mathcal W(\rho,K)$. There is $G=G(\rho,K)$ such that for $0<\varepsilon\le1$, $\hbar>0$, $s_0,s_1\ge0$ and $a\in L^2(\Hb)$,
\[
\norm{W\,V_\varepsilon(s_1,s_0)a}\le\ue^{G\hbar|s_1-s_0|/(2\varepsilon)}\norm{Wa}.
\]
Consequently, for $\gamma\ge0$ and $X_0^\gamma a\in L^2$,
$\norm{X_0^\gamma V_\varepsilon(s_1,s_0)a}\le\ue^{G_\gamma\hbar|s_1-s_0|/(2\varepsilon)}\norm{X_0^\gamma a}$ with $G_\gamma=G(\gamma,K_\gamma)$.
\end{proposition}
\begin{proof}
Let $v(s)=V_\varepsilon(s,s_0)a$ and $h=W^2$. Since $v'=\tfrac{i\hbar}2\Delta_\varepsilon(s)v$ in $L^2$ and $W$ is bounded,
$\tfrac d{ds}\norm{Wv}^2=-\hbar\,\mathrm{Im}\langle hv,\Delta_\varepsilon(s)v\rangle$. Write
\[
\langle hv,J_\varepsilon\Delta(s)J_\varepsilon v\rangle=\langle hJ_\varepsilon v,\Delta(s)J_\varepsilon v\rangle+\langle[J_\varepsilon,h]v,\Delta(s)J_\varepsilon v\rangle .
\]
For $f=J_\varepsilon v$, integrating by parts as in Lemma~\ref{lem:ell},
\[\mathrm{Im}\langle hf,\Delta(s)f\rangle=-\ue^{-2s}\mathrm{Im}\langle(Xh)f,Xf\rangle-\mathrm{Im}\langle(Yh)f,Yf\rangle,\]
which for $s\ge0$ is at
most $4\rho\norm{Wf}\max_Z\norm{WZf}\le4\rho C_W^2\varepsilon^{-1}\norm{Wv}^2$ in absolute value, by Lemma~\ref{lem:wmoll}(i),(ii). By
Lemma~\ref{lem:wmoll}(iv) and its last claim, the second term is at most
$C_W\varepsilon\,\norm{W\Delta(s)J_\varepsilon v}\norm{Wv}\le3C_W^2\varepsilon^{-1}\norm{Wv}^2$. So
$|\tfrac d{ds}\norm{Wv}^2|\le G\hbar\varepsilon^{-1}\norm{Wv}^2$ with $G:=(4\rho+3)C_W^2$, at all times $s\ge0$, and Gr\"onwall's inequality gives
the first claim. The second follows with $W=W_{\gamma,n}$ and $n\to\infty$.
\end{proof}

\begin{lemma}[the path of the mollified flow]\label{lem:path}
Let $\gamma\ge0$, $0<\varepsilon\le1$, $\hbar>0$ and $s_1\ge0$ with $\hbar s_1\le\varepsilon$. Let $\mathcal S:=\{1,X,Y,X^2,Y^2\}$, let $b'$ be such that
$X_0^\gamma Zb'\in L^2(\Hb)$ for every $Z\in\mathcal S$, and put
$g(s):=V_\varepsilon(s,s_1)b'$ for $0\le s\le s_1$. There is $C=C(\gamma)$ such that:
\begin{itemize}
\item[(a)] $g(s)=b'+J_\varepsilon h(s)$ with $h(s):=\tfrac{i\hbar}2\int_{s_1}^s\Delta(r)J_\varepsilon g(r)\,dr$, where $h$ is Lipschitz from
$[0,s_1]$ to $L^2(X_0^{2\gamma}d\nu)$ and $\norm{X_0^\gamma h(s)}\le C\varepsilon^{-1}\norm{X_0^\gamma b'}$;
\item[(b)] the defect $d(s):=(\Delta(s)-\Delta_\varepsilon(s))g(s)$ satisfies
$\norm{X_0^\gamma d(s)}\le C\varepsilon^{-3}\sum_{Z\in\mathcal S}\norm{X_0^\gamma Zb'}$;
\item[(c)] $g$ is $C^1$ from $[0,s_1]$ to $L^2(X_0^{2\gamma}d\nu)$ with $g'=\tfrac{i\hbar}2\Delta_\varepsilon(s)g$, and $Zg$ is
continuous from $[0,s_1]$ to $L^2(X_0^{2\gamma}d\nu)$ for every $Z\in\mathcal S$.
\end{itemize}
\end{lemma}
\begin{proof}
The bounds of Lemma~\ref{lem:wmoll}(i)--(iii) and of Proposition~\ref{prop:ws} hold at the weight $X_0^\gamma$, with $C_W=C_W(\gamma,K_\gamma)$
and $G_\gamma$.
(a) Since $\partial_sg=\tfrac{i\hbar}2J_\varepsilon\Delta(s)J_\varepsilon g$ and $g(s_1)=b'$, the formula for $g$ holds. For $0\le r\le s_1$,
$\hbar|r-s_1|\le\varepsilon$, so $\norm{X_0^\gamma g(r)}\le\ue^{G_\gamma/2}\norm{X_0^\gamma b'}$ and
$\norm{X_0^\gamma\Delta(r)J_\varepsilon g(r)}\le3C_W\varepsilon^{-2}\ue^{G_\gamma/2}\norm{X_0^\gamma b'}$. Hence $h$ is Lipschitz, and
$\norm{X_0^\gamma h(s)}\le\tfrac12\hbar|s-s_1|\cdot3C_W\varepsilon^{-2}\ue^{G_\gamma/2}\norm{X_0^\gamma b'}\le\tfrac32C_W\ue^{G_\gamma/2}\varepsilon^{-1}\norm{X_0^\gamma b'}$.
(b) By (a),
$d(s)=\Delta(s)b'-J_\varepsilon\Delta(s)J_\varepsilon b'+\Delta(s)J_\varepsilon h(s)-J_\varepsilon\Delta(s)J_\varepsilon J_\varepsilon h(s)$. The first term is
bounded by the words $X^2b'$, $Y^2b'$, $Yb'$; the others by $3C_W^2\varepsilon^{-2}\norm{X_0^\gamma b'}$ and
$(3C_W+3C_W^3)\varepsilon^{-2}\norm{X_0^\gamma h(s)}$. (c) follows from (a) and Lemma~\ref{lem:wmoll}: $ZJ_\varepsilon$ is bounded on
$L^2(X_0^{2\gamma}d\nu)$ for $Z\in\mathcal S$, and so is $J_\varepsilon\Delta(s)J_\varepsilon$, continuously in $s$.
\end{proof}

\subsection{Remainder bounds}\label{subsec:dd-wb}
The remainder proof needs two estimates with different roles: factorial derivative bounds for its
amplitude, and exponential smallness of the mollifier defect without a spatial weight.

\begin{lemma}[derivative bounds for the remainder amplitude]\label{lem:wb}
Assume $0<\hbar\le1$ and $0<\alpha\le\alpha_0\le\min(1,\tau_0)$. Let $Z_S$ be a word in $X,Y$ of length at most $n$ (the empty word is the identity). Uniformly in $s\in\Sigma_{K+1}(t)$:
\begin{itemize}
\item[\textup{(i)}]
\[
\norm{X_0^\beta Z_Sb'(s)}\le B_\beta C_4^{N+K+n+1}(2N)!\,(n+2K+2)!\,\alpha^{-(n+2K+2)}(\alpha\hbar)^{-2N}\pwNorm a
\]
\item[\textup{(ii)}] for $0\le\alpha_1<\alpha_0'\le\alpha_0$ with $\alpha_0'-\alpha_1\ge\alpha/4$:
\[
\norm{\ue^{\alpha_1\Lambda}b'(s)}\le C_4^{N+K+1}(2N)!\,(2K+2)!\,\alpha^{-(2K+2)}(\alpha\hbar)^{-2N}\norm{\ue^{\alpha_0'\Lambda}a}.
\]
\end{itemize}
\end{lemma}

\begin{proof}
Expand $b'$ by choosing one term from each factor. Each $\Delta(s_j)$ contributes $\ue^{-2s_j}X^2$, $Y^2$ or $-Y$,
with coefficients of modulus at most $1$ and length at most $2$. Each $\Delta_\varphi(s_1)$ contributes
$\hbar^{-2}(-1+i\hbar)$ (length $0$, coefficient of modulus at most $2\hbar^{-2}$), $-2i\hbar^{-1}Y$
(length $1$, coefficient of modulus $2\hbar^{-1}$), or one of the three terms of length at most $2$ with coefficient
of modulus at most $1$.
Suppose $n_0,n_1,n_2$ of the $N$ $\Delta_\varphi$-factors choose these three kinds. The coefficient is $\le2^{n_0+n_1}\hbar^{-2n_0-n_1}$, and the word length is $\le L:=n+q+2K+2$, with $q:=n_1+2n_2\le2N$.
Since $2n_0+n_1=2N-q$, the coefficient is $\le2^N\hbar^{-2N+q}$.
By Lemma~\ref{lem:sc}, and since $m\mapsto(m+1)m!(C_3/\alpha)^m$ is increasing, each term is $\le B_\beta(L+1)L!(C_3/\alpha)^L\pwNorm a$.
Use $L!\le2^Lq!(n+2K+2)!$, and then
\[
\begin{aligned}
\hbar^{-2N+q}q!(C_3/\alpha)^q&=\hbar^{-2N}q!(C_3\hbar/\alpha)^q\\
&\le\hbar^{-2N}(2N)!\max(1,C_3\hbar/\alpha)^{2N}\le(2N)!\,C_3^{2N}(\alpha\hbar)^{-2N},
\end{aligned}
\]
using $\alpha,\hbar\le1\le C_3$.
Summed over all choices, the factors $2^{n_0+n_1}$ give at most $7^N3^{K+1}$: each $\Delta_\varphi$-factor has two choices
that carry the factor $2$ (lengths $0$ and $1$) and three that carry the factor $1$, so it contributes $2+2+1+1+1=7$, and each
$\Delta(s_j)$ has three choices. Collecting the powers of $2$ and $7$ into $C_4$ gives (i).
Part (ii) follows from the same count, using Lemma~\ref{lem:wc} (gap $\ge\alpha/4$) in place of Lemma~\ref{lem:sc}, and $n=0$. \qedhere
\end{proof}

\begin{lemma}[analytic mollifier defect]\label{lem:unw}
There is a constant $C_0$ such that for $0<\varepsilon\le1$, $0<r\le1$, $s\ge0$ and $v\in\operatorname{Dom}(\ue^{r\Lambda})$,
\[
\norm{(\Delta_\varepsilon(s)-\Delta(s))v}\le C_0\,r^{-4}\,\ue^{-r/(8\varepsilon)}\,\norm{\ue^{r\Lambda}v}.
\]
\end{lemma}

\begin{proof}
\emph{The symbol.} For $\lambda\ge0$: $(1+\lambda^2)(1-j_\varepsilon(\lambda))\le5\,r^{-2}\ue^{-r/(4\varepsilon)}\ue^{r\lambda}$. Indeed
$0\le1-j_\varepsilon(\lambda)=\int g(\lambda-\mu)(1-\chi(\varepsilon\mu))\,d\mu\le\int_{|\mu|\ge1/\varepsilon}g(\lambda-\mu)\,d\mu$.

If $\lambda\le1/(2\varepsilon)$, then $|\lambda-\mu|\ge1/(2\varepsilon)$ on the range of integration, so
  $1-j_\varepsilon(\lambda)\le2\ue^{-1/(8\varepsilon^2)}$. With $u=1/\varepsilon\ge1$ the left side is at most
  $2(1+u^2/4)\ue^{-u^2/8}\le5\ue^{-u/4}\le5\ue^{-r/(4\varepsilon)}$, because $2(1+u^2/4)\ue^{-(u^2-2u)/8}<5$ for $u\ge1$.

If $\lambda>1/(2\varepsilon)$, then $1-j_\varepsilon\le1$ and $(1+\lambda^2)\ue^{-r\lambda/2}\le1+16/(\ue^2r^2)\le4r^{-2}$, so
  $1+\lambda^2\le4r^{-2}\ue^{r\lambda}\ue^{-r\lambda/2}\le4r^{-2}\ue^{r\lambda}\ue^{-r/(4\varepsilon)}$.

\emph{The splitting.} $\Delta_\varepsilon(s)-\Delta(s)=-J_\varepsilon\Delta(s)(1-J_\varepsilon)-(1-J_\varepsilon)\Delta(s)$.

\emph{First term.} With $v_1:=(1-J_\varepsilon)v$, Lemma~\ref{lem:ell}, $\norm{J_\varepsilon}\le1$ and the symbol bound give
\[
\begin{aligned}
\norm{J_\varepsilon\Delta(s)v_1}&\le2(\norm{\Delta v_1}+\norm{v_1})
\le4\norm{(1+\Lambda^2)(1-j_\varepsilon(\Lambda))v}\\
&\le20\,r^{-2}\ue^{-r/(4\varepsilon)}\norm{\ue^{r\Lambda}v}.
\end{aligned}
\]

\emph{Second term.} By Lemma~\ref{lem:wc} with $\tau_0=1$ (words of length $1$ and $2$, radii $r/2<r$),
$\norm{\ue^{(r/2)\Lambda}\Delta(s)v}\le3(1+4C_1/r)^2\norm{\ue^{r\Lambda}v}$. By the symbol bound with $r/2$ in place of $r$,
$\norm{(1-J_\varepsilon)w}\le20\,r^{-2}\ue^{-r/(8\varepsilon)}\norm{\ue^{(r/2)\Lambda}w}$. Take $w=\Delta(s)v$.

Adding the two terms gives the claim with $C_0=20+60(1+4C_1)^2$. \qedhere
\end{proof}
